\documentclass[10pt,landscape]{article}
\usepackage[utf8]{inputenc}
\usepackage[landscape,margin=0.5in]{geometry}
\usepackage{multicol}
\usepackage{amsmath}
\usepackage{amssymb}
\usepackage{xcolor}
\usepackage{enumitem}
\usepackage{titlesec}
\usepackage{array}
\usepackage{booktabs}

% Reduce spacing
\setlength{\parindent}{0pt}
\setlength{\parskip}{2pt}
\setlist{nosep,leftmargin=*}

% Section formatting
\titleformat{\section}{\large\bfseries\color{blue!70!black}}{\thesection}{0.5em}{}
\titleformat{\subsection}{\normalsize\bfseries\color{blue!50!black}}{\thesubsection}{0.5em}{}
\titlespacing*{\section}{0pt}{4pt}{2pt}
\titlespacing*{\subsection}{0pt}{3pt}{1pt}

% Custom commands
\newcommand{\slide}[1]{\section*{Slide #1}}
\newcommand{\pro}{\textcolor{green!60!black}{$\checkmark$}}
\newcommand{\con}{\textcolor{red}{$\times$}}

\pagestyle{empty}

\begin{document}

\begin{center}
\Large\textbf{LSCM Presentation Cheatsheet}\\
\normalsize Least Squares Conformal Maps (SIGGRAPH 2002)
\end{center}

\begin{multicols}{3}
\small

\slide{1: Title}
\textbf{Paper:} LSCM for Automatic Texture Atlas Generation\\
\textbf{Authors:} Lévy, Petitjean, Ray, Maillot\\
\textbf{Presenter:} Tanner Muscarella

\slide{2: The Problem}
\textbf{Goal:} Paint high-res details on 3D models

\textbf{Challenge:}
\begin{itemize}
\item Map 3D surface $\rightarrow$ 2D texture
\item Cut into charts (disk topology)
\item Parameterize each chart
\end{itemize}

\textbf{Prior methods problems:}
\begin{itemize}
\item Require \textbf{convex boundaries}
\item Generate \textbf{many small charts}
\item Seams, artifacts, hard to paint
\end{itemize}

\textcolor{red}{\textbf{Key:}} Need fewer, larger charts with natural boundaries

\slide{3: Conformal Mappings}
\textbf{Definition:} Preserves angles locally

\textbf{Cauchy-Riemann condition:}
\[
\frac{\partial U}{\partial x} + i \frac{\partial U}{\partial y} = 0
\]
where $U = u + iv$ (complex)

\textbf{Meaning:}
\begin{itemize}
\item Iso-$u$ and iso-$v$ tangents: orthogonal, same length
\item Maps circles to circles locally
\end{itemize}

\textbf{Trade-off:} Cannot preserve BOTH angles AND areas (Theorema Egregium)

\textcolor{red}{\textbf{Key insight:}} Use complex numbers to couple $u$ and $v$ directly!

\slide{4: LSCM Method}
\textbf{Key Idea:} Minimize Cauchy-Riemann violation (least-squares)

\textbf{Per-Triangle:}
\[
C(T) = \left| \frac{\partial U}{\partial x} + i \frac{\partial U}{\partial y} \right|^2 A_T
\]

\textbf{Global:}
\[
C(\mathbf{U}) = \sum_{T} C(T) = \mathbf{U}^* \mathbf{C} \mathbf{U}
\]
where $\mathbf{C} = \mathbf{M}^* \mathbf{M}$ (Hermitian Gram matrix)

\textbf{Result:}
\begin{itemize}
\item Quadratic $\Rightarrow$ \textbf{sparse linear system}
\item Pin only \textbf{2 vertices} $\Rightarrow$ unique solution
\item Solve with Conjugate Gradient (fast!)
\end{itemize}

\textbf{Algorithm:}
\begin{enumerate}
\item Build sparse matrix $\mathbf{M}$ (one row per triangle)
\item Pin 2 boundary vertices
\item Solve $(\mathbf{A}^\top \mathbf{A})\mathbf{x} = \mathbf{A}^\top \mathbf{b}$
\item Extract UV coordinates
\end{enumerate}

\slide{5: Why LSCM is Better}
\textbf{Four Guarantees:}

\begin{enumerate}
\item \textbf{Unique Minimum}\\
No local minima (vs MIPS)

\item \textbf{Free Boundary}\\
Boundary moves freely\\
No convex polygon needed\\
Complex borders OK!

\item \textbf{No Triangle Flips}\\
Orientation preserved\\
$\det(J) > 0$ always

\item \textbf{Resolution Independent}\\
Unchanged when mesh refined\\
Good for LOD
\end{enumerate}

\textbf{Practical:} Fast sparse solve vs expensive non-linear optimization

\slide{6: Pipeline}
\textbf{Three Steps:}

\textbf{1. Segmentation}
\begin{itemize}
\item Detect features (high curvature)
\item Region growing from seeds
\item Charts meet at features
\end{itemize}

\textbf{2. Parameterization}
\begin{itemize}
\item Apply LSCM to each chart
\item Pin 2 boundary vertices
\item Solve sparse linear system
\end{itemize}

\textbf{3. Packing}
\begin{itemize}
\item ``Tetris'' algorithm
\item Handles complex borders
\item Better than bounding rectangles
\end{itemize}

\textcolor{red}{\textbf{Advantage:}} Free boundary $\Rightarrow$ larger charts $\Rightarrow$ fewer discontinuities

\columnbreak

\slide{7: Results - Large Charts}
\textbf{Horse Body:}
\begin{itemize}
\item \textbf{72,438 triangles} in single chart!
\item Complex, non-convex boundary
\item Free boundary extrapolated
\item Robust despite large triangles
\end{itemize}

\textbf{Comparison (Cow Head):}

\begin{tabular}{lcc}
\toprule
Method & Before & After \\
\midrule
Harmonic Maps & 3.2 & 1.65 \\
\textbf{LSCM} & \textbf{3.5} & \textbf{1.52} \\
\bottomrule
\end{tabular}

Similar quality, LSCM enables larger charts

\slide{8: More Results}
\textbf{Statistics:}
\begin{itemize}
\item Bunny: 23 charts (69K faces)
\item Horse: 44 charts (97K faces)
\item Natural boundaries
\item Low angle distortion
\end{itemize}

\textbf{Benefits:}
\begin{itemize}
\item Fewer charts vs prior methods
\item Follow geometric features
\item Easier to paint large areas
\item Good for procedural textures
\end{itemize}

\slide{9: Comparison}
\subsection*{1. Harmonic Maps (1994)}
\pro\ Fast, similar quality\\
\con\ Requires convex boundary\\
\con\ Indirect coupling ($u$, $v$ separate)\\
\textbf{Limitation:} Many small charts

\subsection*{2. MIPS (2000)}
\pro\ Free boundary\\
\con\ Non-linear (expensive, local minima)\\
\con\ No guarantees\\
\textbf{LSCM:} Same flexibility + guarantees + faster

\slide{10: Impact \& Limitations}
\textbf{Impact (20+ years):}
\begin{itemize}
\item Still widely used in production
\item Foundation for ARAP, multi-chart methods
\item In 3D paint, game engines, animation
\end{itemize}

\textbf{Limitations:}
\begin{itemize}
\item \textbf{Area distortion:} Conformal $\Rightarrow$ area variation (fixable via post-processing)
\item \textbf{Overlaps possible:} If boundary self-intersects (rare, auto-detected)
\item \textbf{Needs cutting:} Charts must be disk topology
\end{itemize}

\slide{11: Takeaways}
\textbf{Three Contributions:}

\textbf{1. LSCM Parameterization}
\begin{itemize}
\item Fast, robust, flexible
\end{itemize}

\textbf{2. Complex Formulation}
\begin{itemize}
\item Direct $u$-$v$ coupling
\item Elegant Cauchy-Riemann
\item Quadratic minimization
\end{itemize}

\textbf{3. Complete Pipeline}
\begin{itemize}
\item Feature segmentation
\item LSCM parameterization
\item Efficient packing
\end{itemize}

\textcolor{red}{\textbf{Bottom line:}} Large charts with natural boundaries $\Rightarrow$ better texture atlases

\columnbreak

\section*{Key Talking Points}

\subsection*{Opening (Slide 2)}
\textit{``The problem this paper solves is creating texture atlases for 3D paint systems. Prior methods required convex boundaries, forcing many small charts with visible seams.''}

\subsection*{Core Innovation (Slide 4)}
\textit{``The key insight is using complex numbers to directly couple $u$ and $v$ coordinates. This transforms the conformal mapping into a simple quadratic minimization solved with a sparse linear system.''}

\subsection*{Why It Matters (Slide 5)}
\textit{``Unlike previous methods, LSCM provides mathematical guarantees: unique solution, no triangle flips, and free boundaries. This enables much larger charts.''}

\subsection*{Impact (Slide 10)}
\textit{``Twenty years later, this method is still widely used in production and has influenced many follow-up papers.''}

\section*{Common Q\&A}

\textbf{Q: Why complex numbers?}\\
A: Elegantly couple $u$ and $v$ in one equation. Cauchy-Riemann naturally expresses conformality. Multiplication by $i$ = 90° rotation.

\textbf{Q: What if boundary self-intersects?}\\
A: LSCM guarantees no triangle flips, but boundary can self-intersect in UV, creating overlaps. Auto-detected, rare. If happens, subdivide chart.

\textbf{Q: Compare to modern methods?}\\
A: Still competitive. ARAP builds on LSCM but optimizes different criteria (less area distortion).

\textbf{Q: Computational cost?}\\
A: $O(n)$ matrix construction, $O(n \log n)$ to $O(n^2)$ for CG solve. Much faster than non-linear methods.

\textbf{Q: Why only 2 pinned vertices?}\\
A: Need $\geq 2$ to remove translation/rotation ambiguity. Exactly 2 is optimal---allows maximum conformality. More pins over-constrain.

\section*{Technical Details}

\subsection*{Matrix M construction:}
One row per triangle, one column per vertex\\
$M[i,j] = W_j / \sqrt{2 \cdot \text{Area}_i}$ if vertex $j$ in triangle $i$

\subsection*{Weight computation:}
For triangle with local coords $(x_1,y_1), (x_2,y_2), (x_3,y_3)$:
\begin{align*}
W_1 &= (x_3-x_2) + i(y_3-y_2)\\
W_2 &= (x_1-x_3) + i(y_1-y_3)\\
W_3 &= (x_2-x_1) + i(y_2-y_1)
\end{align*}
Note: $W_1 + W_2 + W_3 = 0$ always

\subsection*{Stretch metric:}
Texture area / model area ratio. Lower = better. LSCM: 1.52 after optimization (competitive)

\section*{Timing (10 min)}
\begin{tabular}{ll}
Slide 1 & 15 sec\\
Slide 2 & 1 min\\
Slide 3 & 1 min\\
Slide 4 & 1.5 min \textbf{(most important)}\\
Slide 5 & 1 min\\
Slide 6 & 45 sec\\
Slide 7 & 1 min\\
Slide 8 & 45 sec\\
Slide 9 & 1.5 min\\
Slide 10 & 1 min\\
Slide 11 & 45 sec\\
\end{tabular}

\vspace{4pt}
\textbf{Total:} $\sim$10 min (leave time for questions)

\section*{Remember}
\begin{itemize}
\item Speak clearly, not too fast
\item Eye contact with audience
\item Point to figures when discussing
\item Emphasize \textbf{free boundary} (main contribution)
\item If short on time, abbreviate slides 8-9
\item Have paper open for reference
\end{itemize}

\end{multicols}

\end{document}
